\documentclass[11pt,aspectratio=169]{beamer}

\usetheme{Madrid}
\usecolortheme{default}
\setbeamertemplate{navigation symbols}{}

\usepackage{amsmath,amssymb}
\usepackage{booktabs}

\title{Economic Shocks and Intimate Partner Violence}
\subtitle{Four Papers on Trade, Education, and Gender}
\author{Erika Lynet Salvador}
\date{}

\begin{document}

% --------------------------------------------------
\begin{frame}
\titlepage
\end{frame}

% --------------------------------------------------
\begin{frame}{Big Question}

How do changes in women's economic position affect
intimate partner violence (IPV)?

\medskip

Across the four papers:

\begin{itemize}
    \item Trade shocks change men’s and women’s employment.
    \item Education reform changes women’s human capital.
\end{itemize}

\medskip

The key issue:
Does empowerment reduce violence — or can it increase it?

\end{frame}

% ==================================================
\section{Paper 1: Turkey (Refugees)}

\begin{frame}{Paper 1: Research Question}

\textbf{Erten \& Keskin (JDE, 2021)}

\medskip

Does female employment reduce domestic violence?

\medskip

The authors study Syrian refugee inflows in Turkey,
which reduced employment opportunities for low-skilled native women.

\end{frame}

\begin{frame}{Paper 1: Methodology}

\begin{itemize}
    \item Refugees entered some provinces more than others.
    \item The authors use a distance-based instrument to predict refugee exposure.
    \item Compare IPV before and after refugee inflows.
\end{itemize}

\medskip

Goal: isolate exogenous changes in women’s employment.

\end{frame}

\begin{frame}{Paper 1: Results}

Key findings:

\begin{itemize}
    \item One standard deviation increase in refugee exposure
    reduces physical violence by about 2 percentage points.
    \item This is about a 22\% decline relative to the mean.
    \item Similar declines in sexual and psychological violence.
\end{itemize}

\medskip

Interpretation:
When women’s employment opportunities fall,
violence falls.

This supports a \textbf{rent-extraction theory} of violence.

\end{frame}

\begin{frame}{Paper 1: Why This Matters for Us}

This paper shows:

\begin{itemize}
    \item Labor market shocks can causally change IPV.
    \item Violence may respond to economic incentives.
\end{itemize}

For our project:
It reminds us that shifts in employment conditions
may affect violence through relative income changes.

\end{frame}

% ==================================================
\section{Paper 2: Cambodia (Trade)}

\begin{frame}{Paper 2: Research Question}

\textbf{Erten \& Keskin (REStat, 2024)}

\medskip

How did Cambodia’s WTO accession affect IPV?

\medskip

Trade reduced tariffs unevenly across districts.
This created local labor demand shocks.

\end{frame}

\begin{frame}{Paper 2: Methodology}

\begin{itemize}
    \item District exposure measured using pre-1998 industry shares.
    \item Compare more exposed vs less exposed districts.
    \item Use DHS survey data to measure violence.
\end{itemize}

\medskip

This is a shift-share design.

\end{frame}

\begin{frame}{Paper 2: Labor Market Effects}

In more exposed districts:

\begin{itemize}
    \item Men’s paid employment declines.
    \item Women increase labor force participation.
    \item Many women enter unpaid family work.
\end{itemize}

\end{frame}

\begin{frame}{Paper 2: Violence Effects}

In districts with larger tariff declines:

\begin{itemize}
    \item Physical violence increases by about 0.3 percentage points.
    \item Sexual and psychological violence also increase.
    \item Injury measures increase.
\end{itemize}

\medskip

Although small in absolute terms,
this slowed the overall decline in IPV by about 4\%.

\end{frame}

\begin{frame}{Paper 2: Interpretation}

The authors rule out:

\begin{itemize}
    \item Changes in marriage or fertility
    \item Large consumption declines
    \item Migration effects
\end{itemize}

\medskip

Most consistent explanation:
When male status declines and women begin working,
backlash increases violence.

\medskip

For our project:
Relative employment changes may matter more than total employment.

\end{frame}

% ==================================================
\section{Paper 3: Turkey (Education)}

\begin{frame}{Paper 3: Research Question}

\textbf{Erten \& Keskin (AEJ Applied, 2018)}

\medskip

Does increasing women’s education reduce violence?

\medskip

The authors use Turkey’s 1997 schooling reform,
which increased compulsory education from 5 to 8 years.

\end{frame}

\begin{frame}{Paper 3: Methodology}

\begin{itemize}
    \item Regression discontinuity design.
    \item Compare women born just before vs just after reform cutoff.
\end{itemize}

\medskip

First stage:
\begin{itemize}
    \item Schooling increases by about 0.9 years.
    \item Junior high completion increases by 19 percentage points.
\end{itemize}

\end{frame}

\begin{frame}{Paper 3: Results}

Physical violence does not significantly change.

However:

\begin{itemize}
    \item Psychological violence increases by about 12 percentage points (rural women).
    \item Financial control increases by about 23 percentage points.
\end{itemize}

\medskip

Interpretation:
Education improves women’s economic position,
but may trigger controlling behavior.

\end{frame}

\begin{frame}{Paper 3: Why This Matters}

This paper shows:

\begin{itemize}
    \item Not all forms of violence move the same way.
    \item Education can increase non-physical forms of coercion.
\end{itemize}

For our project:
Measurement matters, i.e., psychological and economic abuse
may respond differently than physical violence.

\end{frame}

% ==================================================
\section{Paper 4: Brazil (Migration)}

\begin{frame}{Paper 4: Why Include This?}

\textbf{Oliveira (JHC, 2016)}

\medskip

Not directly about IPV,
but helps understand women’s outside options.

\medskip

The paper studies how fertility affects women’s migration
and wages.

\end{frame}

\begin{frame}{Paper 4: Methodology and Results}

\begin{itemize}
    \item Uses twins as an instrument for fertility.
    \item An additional child reduces migration by 6\%.
    \item Migration increases probability of working for pay by 5 percentage points.
    \item Migration increases wages significantly.
\end{itemize}

\medskip

Implication:
Children restrict women’s mobility,
which affects economic independence.

\end{frame}

\begin{frame}{Paper 4: Connection to IPV Research}

If fertility restricts migration,
women may have weaker outside options.

\medskip

For our project:
Understanding mobility and labor constraints
is important for interpreting violence outcomes.

\end{frame}

% ==================================================
\section{Overall Lessons}

\begin{frame}{What Do These Papers Teach Us?}

Across countries:

\begin{itemize}
    \item Male job loss can increase violence.
    \item Female employment gains can increase or decrease violence.
    \item Relative status within marriage is crucial.
\end{itemize}

\medskip

Violence responds to economic incentives and social norms.

\end{frame}

\begin{frame}{How This Informs Our Project}

For our research:

\begin{itemize}
    \item We should identify exogenous economic shocks.
    \item We must distinguish prevalence from reporting.
    \item We should measure multiple dimensions of violence.
\end{itemize}

\end{frame}

% --------------------------------------------------
\begin{frame}{Conclusion}

These four papers show that:

\medskip

\begin{center}
Economic change reshapes household power.
\end{center}

\medskip

Violence is not only about poverty, but about relative status and control.

\end{frame}

\end{document}
